% ==============================================================================
% บทที่ 2: การทบทวนวรรณกรรมที่เกี่ยวข้อง (Literature Review)
% โครงการ: การพัฒนาแพลตฟอร์มสำหรับการสนับสนุนการวินิจฉัยรอยโรคทางช่องปาก
%          และผสมผสานการใช้งานระบบผู้เชี่ยวชาญสำหรับทันตแพทย์
% ==============================================================================

\chapter{การทบทวนวรรณกรรมที่เกี่ยวข้อง}
\label{chapter:literature-review}

การตรวจคัดกรองและวินิจฉัยรอยโรคในช่องปาก (Oral Lesion Screening and Diagnosis) โดยเฉพาะการประเมินความเสี่ยงต่อการเกิดโรคมะเร็งช่องปาก (Oral Cancer) ถือเป็นขั้นตอนที่มีความสำคัญอย่างยิ่งต่อการเพิ่มอัตราการรอดชีวิตของผู้ป่วย โดยมะเร็งช่องปากส่วนใหญ่จัดอยู่ในกลุ่มมะเร็งเซลล์สความัส (Oral Squamous Cell Carcinoma: OSCC) ซึ่งมักมีจุดเริ่มต้นมาจากการเปลี่ยนแปลงของเนื้อเยื่อบุผิวที่เป็นรอยโรคก่อนมะเร็ง (Oral Potentially Malignant Disorders: OPMDs) อย่างไรก็ตาม ในทางคลินิกพบว่ารอยโรคระยะแรกเริ่มมักมีลักษณะทางสัณฐานวิทยาใกล้เคียงกับรอยโรคทั่วไปที่ไม่ร้ายแรง ทำให้การแยกแยะโรคต้องพึ่งพาความชำนาญของทันตแพทย์เฉพาะทางด้านเวชศาสตร์ช่องปาก ซึ่งยังมีจำนวนจำกัดและกระจุกตัวอยู่ในโรงพยาบาลระดับศูนย์

จากการศึกษาและต่อยอดโครงงานวิจัยของรุ่นพี่ในปีก่อนหน้าทั้ง 2 รุ่น พบว่า:
\begin{enumerate}
    \item \textbf{โครงงานวิจัยของรุ่นพี่กลุ่มที่ 1 (ปีการศึกษา 2567):} ได้พัฒนาระบบตรวจจับรอยโรคโดยใช้แบบจำลองตรวจจับวัตถุ (Object Detection) ตระกูล YOLO ร่วมกับการประมวลผลข้อมูลอาการของผู้ป่วยด้วยกฎเงื่อนไข (Rule-based) แม้จะสามารถตีกรอบ Bounding Box ระบุตำแหน่งรอยโรคได้ดี แต่การตัดสินใจยังเป็นการวิเคราะห์แบบแยกส่วน และมีข้อจำกัดอย่างมากเมื่อภาพถ่ายไม่สามารถตรวจจับรอยโรคได้ (No Detection) หรือเมื่อผู้ป่วยมีอาการที่ขัดแย้งกัน
    \item \textbf{โครงงานวิจัยของรุ่นพี่กลุ่มที่ 2 (ปีการศึกษา 2568):} ได้พัฒนาอัลกอริทึมการเรียนรู้เชิงลึกแบบหลายโมดาลิตี (Multimodal Deep Learning) โดยใช้โครงข่าย DenseNet-121 วิเคราะห์ภาพถ่ายร่วมกับแบบจำลองภาษาไทย WangchanBERTa วิเคราะห์ข้อความซักประวัติ ซึ่งให้ค่าความแม่นยำสูงถึงร้อยละ 95.5 และทนทานต่อข้อมูลอาการที่ขัดแย้งกันได้ดี แต่มีข้อจำกัดคือเป็นแบบจำลองที่ทำนายเฉพาะภาพรวม (Whole-image classification) ไม่สามารถตีกรอบระบุพิกัดตำแหน่งรอยโรคบนภาพได้
\end{enumerate}

แม้ว่าระบบต้นแบบของรุ่นพี่ทั้งสองกลุ่มจะประสบความสำเร็จอย่างสูงในระดับห้องปฏิบัติการ แต่เมื่อพิจารณาในมุมมองของวิศวกรรมซอฟต์แวร์และการนำไปให้บริการจริงทางการแพทย์ พบว่าระบบเดิมยังคงอยู่ในระดับ ``งานวิจัยต้นแบบ (Research Prototype)'' ที่มีข้อจำกัดด้านโครงสร้างระบบ เช่น สถาปัตยกรรมที่กินทรัพยากรประมวลผลคลาวด์สูง ทำให้มีต้นทุนค่าบริการแพง ขาดระบบการจัดการวงจรชีวิตซอฟต์แวร์และการส่งมอบระบบอัตโนมัติ (CI/CD) ขาดระบบการทดสอบอัตโนมัติ (Automated Testing) ตลอดจนยังไม่มีระบบยืนยันตัวตนและการจัดการสิทธิ์ที่ปลอดภัยตามมาตรฐานพระราชบัญญัติคุ้มครองข้อมูลส่วนบุคคล (PDPA) และยังขาดหน้าต่างการทำงานสำหรับทันตแพทย์ (Doctor Dashboard) ที่มีระบบวงจรข้อมูลย้อนกลับ (Feedback Loop) สำหรับให้แพทย์ตรวจสอบ ยืนยัน หรือแก้ไขผลลัพธ์ของ AI เพื่อเก็บรวบรวมเป็นชุดข้อมูลคุณภาพสูงไปพัฒนาแบบจำลองให้แม่นยำยิ่งขึ้นในอนาคต

ดังนั้น เพื่อเป็นพื้นฐานในการออกแบบและพัฒนาแพลตฟอร์มสนับสนุนการตัดสินใจสำหรับทันตแพทย์ที่พร้อมใช้งานจริง (Production-Ready Platform) และบูรณาการระบบผู้เชี่ยวชาญทั้งสองแนวทางเข้าด้วยกัน เนื้อหาในบทนี้จึงมุ่งเน้นการทบทวนวรรณกรรม สถาปัตยกรรมระบบ วิศวกรรมซอฟต์แวร์ และเทคโนโลยีที่เกี่ยวข้อง โดยแบ่งออกเป็น 7 หัวข้อหลัก ได้แก่:
\begin{enumerate}
    \item สถาปัตยกรรมซอฟต์แวร์และระบบคลาวด์ไมโครเซอร์วิส
    \item เทคโนโลยีส่วนต่อประสานผู้ใช้ แดชบอร์ดแพทย์ และการสื่อสาร
    \item วิศวกรรมซอฟต์แวร์ การทดสอบอัตโนมัติ และความปลอดภัยของข้อมูลสุขภาพ
    \item ภาพรวมสถาปัตยกรรมปัญญาประดิษฐ์และมาตรวัดประสิทธิภาพ
    \item การทบทวนงานวิจัยและปริญญานิพนธ์ที่เกี่ยวข้อง
    \item สรุปแนวทางการออกแบบสถาปัตยกรรมและกรอบแนวคิดของโครงงาน
\end{enumerate}

% ------------------------------------------------------------------------------
% 2.1 แนวคิดทางการแพทย์และพยาธิวิทยาของรอยโรคในช่องปาก
% ------------------------------------------------------------------------------
\section{แนวคิดทางการแพทย์และพยาธิวิทยาของรอยโรคในช่องปาก}
\label{sec:medical_pathology}

\subsection{รอยโรคเป้าหมาย 5 กลุ่ม}
ในการพัฒนาระบบคัดกรอง โครงงานนี้ครอบคลุมกลุ่มรอยโรคสีขาวและแผลในช่องปากที่มีความสำคัญทางคลินิก 5 ประเภท ได้แก่:
\begin{enumerate}
    \item \textbf{แผลเปื่อยหรือแผลร้อนใน (Aphthous Ulcer / Ulcer):} เป็นแผลหลุมตื้น ขอบแผลแดง ตรงกลางเป็นสีขาวเหลือง เกิดจากการอักเสบ ผู้ป่วยมีอาการเจ็บปวดชัดเจนเมื่อสัมผัส มักหายได้เองภายใน 1--2 สัปดาห์ แต่หากแผลเรื้อรังเกิน 3 สัปดาห์โดยไม่หายหรือมีขอบแข็งยกตัวขึ้น จะต้องสงสัยมะเร็งช่องปาก
    \item \textbf{โรคติดเชื้อราแคนดิดาในช่องปาก (Oral Candidiasis):} เกิดจากเชื้อรา \textit{Candida albicans} มีลักษณะเป็นแผ่นฝ้าขาวคล้ายคราบนม \textbf{ลักษณะเด่นคือสามารถใช้ผ้ากอซขูดหรือเช็ดออกได้ (Scrapable plaque)} และเมื่อเช็ดออกจะพบผิวเยื่อบุสีแดงด้านล่าง
    \item \textbf{เส้นขาวแนวกระพุ้งแก้ม (Linea Alba Buccalis):} เป็นการเปลี่ยนแปลงทางกายวิภาคปกติ (Normal variation) เกิดจากแรงกดหรือแรงเสียดสีซ้ำๆ จากแนวสบฟัน เป็นเส้นขาวนูนตามแนวนอนของกระพุ้งแก้ม \textbf{ไม่มีอาการเจ็บหรือแสบใดๆ และขูดไม่ออก} ไม่มีความเสี่ยงต่อการเกิดมะเร็ง
    \item \textbf{รอยโรคแผ่นขาว (Oral Leukoplakia):} เป็นรอยโรคก่อนมะเร็ง (OPMD) ที่สำคัญที่สุด มีลักษณะเป็นแผ่นหรือปื้นสีขาวหนาตัว \textbf{ขูดไม่ออก} สัมพันธ์กับพฤติกรรมการสูบบุหรี่ ดื่มสุรา และเคี้ยวหมากเรื้อรัง มีอัตราการกลายพันธุ์เป็นมะเร็งช่องปากประมาณร้อยละ 1--9
    \item \textbf{โรคไลเคน พลานัสในช่องปาก (Oral Lichen Planus: OLP):} เป็นโรคอักเสบเรื้อรังจากระบบภูมิคุ้มกัน มีลักษณะเป็นเส้นสีขาวสานกันคล้ายร่างแห (Wickham's striae) หรือแผลแดงถลอก ผู้ป่วยมักมีอาการแสบร้อนเมื่อรับประทานอาหารเผ็ด จัดเป็นรอยโรคกลุ่ม OPMD ที่มีโอกาสกลายเป็นมะเร็งได้ประมาณร้อยละ 0.5--2
\end{enumerate}

\begin{table}[htbp]
\centering
\small
\caption{สรุปการเปรียบเทียบลักษณะทางคลินิกและพยาธิวิทยาของรอยโรคเป้าหมายทั้ง 5 ชนิด}
\label{tab:clinical_pathology_summary}
\begin{tabular}{p{2.8cm}p{3.2cm}p{2.6cm}p{4.2cm}}
\toprule
\textbf{ชื่อรอยโรค} & \textbf{ลักษณะทางกายภาพ} & \textbf{การขูดลอก} & \textbf{อาการทางคลินิกและพฤติกรรมเสี่ยง} \\
\midrule
Ulcer (แผลเปื่อย) & แผลหลุมตื้น ตรงกลางขาวเหลือง ขอบแดง & ขูดไม่ออก & เจ็บปวดมากเวลากินอาหารหรือสัมผัส หายใน 2 สัปดาห์ \\
Candidiasis (เชื้อรา) & แผ่นฝ้าขาวคล้ายคราบนมกระจายตัว & เช็ด/ขูดลอกได้ & แสบปาก ผิวใต้รอยแดง อาจพบจุดเลือดซึม \\
Linea Alba & เส้นขาวนูนทอดตามแนวสบฟัน & ขูดไม่ออก & ไม่มีอาการใดๆ, สัมพันธ์กับแรงสบฟัน \\
Leukoplakia (แผ่นขาว) & แผ่นหรือปื้นขาวขอบเขตชัดเจน & ขูดไม่ออก & สัมพันธ์กับสูบบุหรี่ ดื่มสุรา เคี้ยวหมาก (เสี่ยงมะเร็ง) \\
Lichen Planus & เส้นขาวร่างแห หรือแผลแดงถลอก & ขูดไม่ออก & ปวดแสบปวดร้อนเมื่อกินอาหารเผ็ด (เสี่ยงมะเร็ง) \\
\bottomrule
\end{tabular}
\end{table}

% ------------------------------------------------------------------------------
% 2.2 สถาปัตยกรรมซอฟต์แวร์และระบบคลาวด์ไมโครเซอร์วิส
% ------------------------------------------------------------------------------
\section{สถาปัตยกรรมซอฟต์แวร์และระบบคลาวด์ไมโครเซอร์วิส}
\label{sec:software_architecture}

\subsection{สถาปัตยกรรมแบบ Monolith เทียบกับสถาปัตยกรรม Microservice}
ในระบบต้นแบบเดิมของงานวิจัยรุ่นก่อนหน้า มักใช้สถาปัตยกรรมแบบ Monolith ซึ่งรวมทั้งส่วนติดต่อผู้ใช้ การประมวลผลทางธุรกิจ และการโหลดโมเดลปัญญาประดิษฐ์เข้าไว้ในแอปพลิเคชันเดียวกัน ส่งผลให้เกิดข้อจำกัดหลายประการ:
\begin{itemize}
    \item \textbf{การใช้ทรัพยากรที่ไม่คุ้มค่า:} การรันโมเดล AI ขนาดใหญ่ (เช่น PyTorch, CUDA, Transformers) ภายในเว็บเซิร์ฟเวอร์หลัก ทำให้ตัวเซิร์ฟเวอร์ต้องใช้หน่วยความจำ (RAM) และ CPU สูงตลอดเวลา แม้จะไม่มีการส่งภาพเข้ามาวิเคราะห์
    \item \textbf{ความยากในการขยายขนาด (Scalability):} เมื่อมีผู้ใช้งานเข้าดูเว็บไซต์พร้อมกันจำนวนมาก ระบบไม่สามารถขยายเฉพาะเว็บเซิร์ฟเวอร์ได้ แต่ต้องคัดลอกทั้งอินสแตนซ์ที่รวมโมเดล AI ไปด้วย ซึ่งทำให้ค่าใช้จ่ายบนคลาวด์เพิ่มสูงขึ้นอย่างรวดเร็ว
    \item \textbf{ความเสี่ยงของระบบล่ม (Single Point of Failure):} หากกระบวนการอนุมานผลของ AI เกิดข้อผิดพลาดจากหน่วยความจำไม่เพียงพอ (Out of Memory: OOM) จะส่งผลให้ทั้งเว็บเซิร์ฟเวอร์หยุดทำงานทันที
\end{itemize}

ด้วยเหตุนี้ โครงงานนี้จึงปรับเปลี่ยนสถาปัตยกรรมไปสู่ สถาปัตยกรรมไมโครเซอร์วิส (Microservices Architecture) โดยแยกส่วนการทำงานออกเป็น 2 เซอร์วิสหลัก:
\begin{enumerate}
    \item \textbf{API Gateway Service:} รับคำขอจากผู้ใช้งาน ตรวจสอบสิทธิ์ ตรวจสอบความถูกต้องของข้อมูล และส่งต่องานไปยัง AI Engine
    \item \textbf{AI Engine Inference Service:} ทำหน้าที่โหลดโมเดลเข้าสู่หน่วยความจำและรันการอนุมานผลภาพและข้อความประวัติผู้ป่วยโดยเฉพาะ
\end{enumerate}

\subsection{รูปแบบประตูเชื่อมต่อระบบ (API Gateway Pattern) และ FastAPI}
FastAPI ถูกเลือกใช้เป็นเฟรมเวิร์กหลักในการสร้าง API Gateway และ AI Engine ด้วยคุณสมบัติที่โดดเด่น:
\begin{itemize}
    \item \textbf{Asynchronous Programming (ASGI):} ทำงานบนสถาปัตยกรรมแบบ Asynchronous ควบคู่กับเซิร์ฟเวอร์ Uvicorn ทำให้สามารถรองรับคำขอ (Concurrency) ได้พร้อมกันจำนวนมหาศาลโดยไม่เกิดการบล็อก (Non-blocking I/O)
    \item \textbf{Data Validation \& Serialization ด้วย Pydantic:} ทำหน้าที่ตรวจสอบโครงสร้างข้อมูลอินพุต (เช่น ขนาดและรูปแบบไฟล์ภาพ, ข้อความประวัติ) และจัดรูปแบบผลลัพธ์การประเมิน (Unified Assessment Schema) ให้เป็นมาตรฐานเดียวกันก่อนส่งกลับไปยังผู้ใช้
    \item \textbf{Auto-Generated Documentation:} สร้างเอกสารแนะนำการใช้งาน API อัตโนมัติในรูปแบบ OpenAPI (Swagger UI) ช่วยให้การพัฒนาและเชื่อมต่อกับระบบภายนอกเป็นไปอย่างมีมาตรฐาน
\end{itemize}

\subsection{เทคโนโลยีคอนเทนเนอร์ (Containerization: Docker \& Docker Compose)}
Docker ถูกนำมาใช้ในการบรรจุแอปพลิเคชัน โค้ด รันไทม์ และชุดไลบรารีที่จำเป็นทั้งหมด (เช่น PyTorch, OpenCV, Transformers) เข้าเป็นอิมเมจมาตรฐาน (Container Images) เพื่อตัดปัญหา ``ทำงานได้บนเครื่องผู้พัฒนาแต่รันบนเซิร์ฟเวอร์ไม่ได้ (It works on my machine problem)'' นอกจากนี้ยังใช้ Docker Compose ในการจัดการประสานงานระหว่างคอนเทนเนอร์ Gateway และคอนเทนเนอร์ AI Engine ทำให้สามารถจำลองสภาพแวดล้อมเสมือนจริงได้อย่างสมบูรณ์ทั้งในเครื่องทดสอบและบนคลาวด์


% ------------------------------------------------------------------------------
% 2.3 เทคโนโลยีส่วนต่อประสานผู้ใช้ แดชบอร์ดแพทย์ และการสื่อสาร
% ------------------------------------------------------------------------------
\section{เทคโนโลยีส่วนต่อประสานผู้ใช้ แดชบอร์ดแพทย์ และการสื่อสาร}
\label{sec:ui_dashboard_comm}

\subsection{การพัฒนาเว็บแอปพลิเคชันส่วนติดต่อผู้ใช้ (Web Application Frontend)}
ส่วนติดต่อผู้ใช้ถูกออกแบบเพื่อรองรับการใช้งานที่สะดวกและรวดเร็วสำหรับทันตแพทย์ โดยใช้เทคโนโลยีเว็บสมัยใหม่:
\begin{itemize}
    \item \textbf{Single-Page Application (SPA):} ช่วยให้การตอบสนองต่อการคลิกและเปลี่ยนหน้าเป็นไปอย่างรวดเร็วโดยไม่ต้องรีโหลดหน้าเว็บใหม่ทั้งหมด
    \item \textbf{Tailwind CSS:} เฟรมเวิร์กการจัดแต่งรูปแบบการแสดงผลแบบ Utility-first ที่ช่วยให้การจัดวางองค์ประกอบมีความยืดหยุ่นสูง สวยงาม รองรับการแสดงผลทุกขนาดหน้าจอ (Responsive Web Design) ทั้งบนคอมพิวเตอร์ แท็บเล็ต และสมาร์ตโฟน
    \item \textbf{Form \& File Upload Handling:} มีระบบตรวจสอบความสมบูรณ์ของภาพถ่ายและแบบฟอร์มซักประวัติก่อนส่งข้อมูล เพื่อลดความผิดพลาดที่อาจเกิดขึ้นจากฝั่งผู้ใช้งาน
\end{itemize}

\subsection{แดชบอร์ดสนับสนุนการตัดสินใจของทันตแพทย์และระบบวงจรข้อมูลย้อนกลับ (Doctor Clinical Dashboard \& Feedback Loop)}
หนึ่งในนวัตกรรมสำคัญของแพลตฟอร์มนี้คือ **ระบบวงจรข้อมูลย้อนกลับ (Doctor Feedback Loop)** ซึ่งถูกสร้างขึ้นบนแดชบอร์ดสำหรับทันตแพทย์โดยเฉพาะ:
\begin{enumerate}
    \item \textbf{Clinical Decision Support:} ระบบจะแสดงผลการวิเคราะห์ภาพถ่ายคู่กับแผนที่ความร้อน Grad-CAM พร้อมแสดงกรอบ Bounding Box และคะแนนระดับความเสี่ยงมะเร็ง (Risk Score) เพื่อให้ทันตแพทย์ใช้ประกอบการวินิจฉัย
    \item \textbf{Doctor Verification \& Relabeling:} ทันตแพทย์ผู้เชี่ยวชาญสามารถกดยืนยันผลการทำนายของ AI หรือทำการแก้ไขชนิดรอยโรค (Relabeling) ได้หากพบว่าผลการทำนายของ AI ยังไม่ถูกต้อง
    \item \textbf{Ground Truth Data Curation:} ข้อมูลที่ได้รับการยืนยันหรือแก้ไขจากแพทย์จะถูกจัดเก็บเข้าสู่ฐานข้อมูลคุณภาพสูง (Curated Ground Truth Dataset) โดยอัตโนมัติ ซึ่งทำหน้าที่เป็นขุมข้อมูลสำคัญสำหรับการนำไปฝึกฝนและปรับปรุงโมเดล AI ในรอบถัดไป (Continuous Learning) ทำให้ระบบมีความแม่นยำเพิ่มขึ้นอย่างยั่งยืน
\end{enumerate}

% ------------------------------------------------------------------------------
% 2.4 วิศวกรรมซอฟต์แวร์ การทดสอบอัตโนมัติ และความปลอดภัยของข้อมูลสุขภาพ
% ------------------------------------------------------------------------------
\section{วิศวกรรมซอฟต์แวร์ การทดสอบอัตโนมัติ และความปลอดภัยของข้อมูลสุขภาพ}
\label{sec:software_engineering_pdpa}

\subsection{การผสานรวมและการส่งมอบระบบอย่างต่อเนื่อง (CI/CD Pipelines)}
เพื่อให้การดูแลรักษาและการพัฒนาฟีเจอร์ใหม่เป็นไปอย่างรวดเร็วและปลอดภัย โครงงานนี้นำแนวคิด DevOps และกระบวนการ **CI/CD (Continuous Integration / Continuous Deployment)** มาประยุกต์ใช้ผ่าน GitHub Actions:
\begin{itemize}
    \item \textbf{Continuous Integration (CI):} ทุกครั้งที่มีการแก้ไขและพุชโค้ดเข้าสู่คลังจัดเก็บ ระบบจะทำการรันการตรวจสอบรูปแบบโค้ด (Linting) และรันชุดทดสอบอัตโนมัติ (Automated Testing) ทันที
    \item \textbf{Continuous Deployment (CD):} เมื่อผ่านการทดสอบครบถ้วน ระบบจะทำการสร้าง Docker Image และส่งมอบไปยังสภาพแวดล้อมคลาวด์ (Cloud Run) โดยอัตโนมัติ ช่วยลดความผิดพลาดจากการติดตั้งแบบ Manual
\end{itemize}

\subsection{การทดสอบระบบอัตโนมัติ (Automated Testing)}
การสร้างความเชื่อถือได้ของซอฟต์แวร์ทางการแพทย์จำเป็นต้องมีระบบการทดสอบที่ครอบคลุม:
\begin{enumerate}
    \item \textbf{Unit Testing \& Integration Testing:} ทดสอบฟังก์ชันการทำงานระดับย่อยและการเชื่อมต่อระหว่าง API Gateway และ AI Engine โดยใช้ Pytest และ HTTPX Mocking
    \item \textbf{Load \& Stress Testing:} การจำลองการส่งคำขอเข้ามาพร้อมกันจำนวนมากเพื่อวัดระยะเวลาการตอบสนอง (Response Time / Latency) และตรวจสอบขีดจำกัดของระบบในการรองรับการใช้งานพร้อมกันของคลินิกหลายแห่ง
    \item \textbf{Security \& Vulnerability Scanning:} ตรวจสอบช่องโหว่ความปลอดภัยของไลบรารีและแพ็กเกจที่ใช้งาน เพื่อป้องกันการถูกโจมตีทางไซเบอร์
\end{enumerate}

\subsection{ระบบบันทึกประวัติและการตรวจจับข้อผิดพลาดส่วนกลาง (Centralized Logging)}
ระบบติดตั้งกลไก Centralized Logging เพื่อรวบรวมบันทึกเหตุการณ์ (Logs) จากทั้ง Gateway และ AI Engine เข้าสู่ศูนย์กลาง ช่วยให้ผู้ดูแลระบบสามารถติดตามข้อผิดพลาด (Error Tracking), ตรวจสอบความผิดปกติของการประมวลผลโมเดล และวิเคราะห์สาเหตุของปัญหาได้อย่างรวดเร็วแบบเรียลไทม์

\subsection{ความปลอดภัยและการคุ้มครองข้อมูลส่วนบุคคล (PDPA Compliance for Patient Records)}
ข้อมูลภาพถ่ายช่องปากและประวัติการรักษาจัดเป็นข้อมูลส่วนบุคคลที่มีความอ่อนไหวสูง (Sensitive Personal Data) ระบบจึงได้รับการออกแบบให้สอดคล้องตาม **พระราชบัญญัติคุ้มครองข้อมูลส่วนบุคคล พ.ศ. 2562 (PDPA)**:
\begin{itemize}
    \item \textbf{Authentication \& Role-Based Access Control (RBAC):} มีระบบยืนยันตัวตนของผู้ใช้งาน และแบ่งระดับสิทธิ์การเข้าถึงอย่างชัดเจน เช่น สิทธิ์สำหรับทันตแพทย์ทั่วไปในการส่งตรวจ และสิทธิ์สำหรับทันตแพทย์ผู้เชี่ยวชาญในการเข้าถึงหน้าแดชบอร์ดตรวจสอบและแก้ไขผล
    \item \textbf{Data Anonymization \& Pseudonymization:} ข้อมูลภาพถ่ายและประวัติที่นำมาประมวลผลในระบบ AI จะถูกตัดข้อมูลที่สามารถระบุตัวตนบุคคลได้ (เช่น ชื่อ-นามสกุล หมายเลขบัตรประชาชน) ออกจากชุดข้อมูลประมวลผล
    \item \textbf{Data Encryption:} ข้อมูลทั้งหมดจะถูกเข้ารหัสขณะส่งผ่านเครือข่ายด้วยโปรโตคอล HTTPS (TLS 1.3) และเข้ารหัสข้อมูลขณะจัดเก็บในฐานข้อมูล (Encryption at Rest)
\end{itemize}

% ------------------------------------------------------------------------------
% 2.5 ภาพรวมสถาปัตยกรรมปัญญาประดิษฐ์และมาตรวัดประสิทธิภาพ
% ------------------------------------------------------------------------------
\section{ภาพรวมสถาปัตยกรรมปัญญาประดิษฐ์และมาตรวัดประสิทธิภาพ}
\label{sec:ai_architecture_overview}

\subsection{ภาพรวมแบบจำลองตรวจจับวัตถุ YOLO (You Only Look Once)}
สถาปัตยกรรม YOLO (YOLOv7 และ YOLOv8) ถูกใช้ในระบบเพื่อทำหน้าที่ตรวจจับตำแหน่งและตีกรอบ Bounding Box ระบุพิกัดรอยโรคบนภาพถ่าย โดยแบบจำลองจะรับภาพถ่ายช่องปากเข้ามา สกัดคุณลักษณะหลายระดับผ่าน Backbone และ Neck จากนั้น Detection Head จะทำนายพิกัดกรอบ $[x, y, w, h]$ และค่าความมั่นใจ (Confidence Score) พร้อมทั้งใช้กระบวนการ Non-Maximum Suppression (NMS) เพื่อคัดกรองกรอบที่ซ้ำซ้อนกันออกไป ทำให้แพทย์เห็นตำแหน่งที่ชัดเจนของแผลบนภาพ

\subsection{ภาพรวมแบบจำลองการเรียนรู้เชิงลึกแบบหลายโมดาลิตี (Multimodal AI)}
เพื่อประเมินความเสี่ยงต่อการเกิดมะเร็งช่องปากอย่างแม่นยำ ระบบได้ผสานข้อมูลสองรูปแบบ:
\begin{enumerate}
    \item \textbf{โมดาลิตีภาพถ่าย (DenseNet-121):} ทำหน้าที่สกัดคุณลักษณะทางกายภาพของรอยโรค เช่น สี ขอบเขต และพื้นผิว ผ่าน Dense Block ที่มีการเชื่อมต่อแบบหนาแน่น ส่งผลให้ได้เวกเตอร์คุณลักษณะขนาด $1 \times 1024$ มิติ
    \item \textbf{โมดาลิตีข้อความประวัติผู้ป่วย (WangchanBERTa):} ทำหน้าที่ประมวลผลข้อความอาการภาษาไทย เช่น ประวัติการสูบบุหรี่ การดื่มสุรา หรืออาการเจ็บแสบ ผ่านกลไก Self-Attention ส่งผลให้ได้เวกเตอร์ตัวแทนบริบทภาษาไทยขนาด $1 \times 768$ มิติ
    \item \textbf{Late Decision Fusion:} นำคะแนนผลลัพธ์ของโมเดลภาพและโมเดลข้อความมารวมกันด้วยการถ่วงน้ำหนัก โดยให้น้ำหนักภาพร้อยละ 60 และข้อความร้อยละ 40 ซึ่งพิสูจน์แล้วว่าช่วยให้ระบบทนทานต่อข้อมูลอาการที่ขัดแย้งกับภาพได้อย่างดีเยี่ยม
\end{enumerate}
% ------------------------------------------------------------------------------
% 2.6 การทบทวนงานวิจัยและปริญญานิพนธ์ที่เกี่ยวข้อง
% ------------------------------------------------------------------------------
\section{การทบทวนงานวิจัยและปริญญานิพนธ์ที่เกี่ยวข้อง}
\label{sec:literature_review_works}

\subsection{ปริญญานิพนธ์ของรุ่นพี่กลุ่มที่ 1 (ปีการศึกษา 2567)}
\textbf{เปี่ยมเดช ยศเงินทองฟู และ ปุณณวิช สิทธิเสือ (2567) \cite{piamdech2024oral}} ได้พัฒนาปริญญานิพนธ์เรื่อง \textit{``การพัฒนาระบบเว็บผู้เชี่ยวชาญและส่วนต่อประสานแอปพลิเคชันไลน์เพื่อการประเมินความเสี่ยงมะเร็งในช่องปากจากรูปแผลและประวัติผู้ป่วย''} โดยใช้แบบจำลอง YOLO (YOLOv8, YOLOv9, YOLOv10) เพื่อตรวจจับและจำแนกรอยโรคในช่องปาก 5 ชนิด จากชุดข้อมูลภาพถ่ายของคณะทันตแพทยศาสตร์ มหาวิทยาลัยมหิดล จำนวน 1,579 ภาพ ควบคู่กับการนำประวัติผู้ป่วยมาประเมินร่วมผ่านกฎเงื่อนไข (Rule-based) บนระบบคลาวด์ GCP และ LINE API
\begin{itemize}
    \item \textbf{ข้อค้นพบและจุดเด่น:} แบบจำลอง YOLOv8 ให้ความถูกต้องสูงสุดที่ร้อยละ 86 เมื่อนำข้อมูลประวัติผู้ป่วยมาประเมินร่วม และมีจุดเด่นคือสามารถวาดกรอบ Bounding Box ระบุตำแหน่งรอยโรคบนภาพได้ชัดเจน
    \item \textbf{ข้อจำกัดและปัญหาที่พบ:} การทำงานระหว่างการตรวจจับภาพและกฎเงื่อนไขอาการเป็นการประมวลผลแบบแยกส่วน (Decoupled) หากภาพถ่ายไม่ชัดจนโมเดลเกิดภาวะ No Detection ระบบจะไม่สามารถใช้อาการของผู้ป่วยมาช่วยวิเคราะห์ต่อได้เลย นอกจากนี้ การตัดสินใจด้วยกฎเงื่อนไขยังไม่สามารถรับมือกับข้อมูลอาการที่มีความขัดแย้งกันได้ดีพอ
\end{itemize}

\subsection{ปริญญานิพนธ์ของรุ่นพี่กลุ่มที่ 2 (ปีการศึกษา 2568)}
\textbf{รตา อุณหะ และ วรดา สาระคนธ์ (2568) \cite{unnahasarakhon2025multimodal}} ได้พัฒนาปริญญานิพนธ์เรื่อง \textit{``การพัฒนาอัลกอริทึมสำหรับระบบผู้เชี่ยวชาญวิเคราะห์รอยโรคทางช่องปากผ่าน Web และ Line API''} โดยมุ่งเน้นการแก้ปัญหาการประมวลผลแบบแยกส่วนของรุ่นพี่กลุ่มที่ 1 ด้วยการพัฒนาแบบจำลองการเรียนรู้เชิงลึกแบบหลายโมดาลิตี (Multimodal Deep Learning) ที่ผสานภาพถ่าย (DenseNet-121) เข้ากับข้อความประวัติผู้ป่วยภาษาไทย (WangchanBERTa) และเปรียบเทียบวิธีการรวมข้อมูลทั้ง Early Fusion และ Late Fusion
\begin{itemize}
    \item \textbf{ข้อค้นพบและจุดเด่น:} วิธีการ Late Fusion แบบกำหนดน้ำหนักภาพต่อข้อความ 60:40 ให้ประสิทธิภาพสูงสุด โดยมีค่าความแม่นยำถึงร้อยละ 95.5 และยังคงรักษาระดับความแม่นยำได้ที่ร้อยละ 82.5 ในสถานการณ์จำลองที่ข้อมูลอาการไม่ตรงกับภาพ นอกจากนี้ยังมีภาพความร้อน Grad-CAM อธิบายการตัดสินใจ
    \item \textbf{ข้อจำกัดและปัญหาที่พบ:} แบบจำลองเป็นแบบการจำแนกภาพทั้งภาพ (Whole-image classification) ทำให้ไม่สามารถตีกรอบ Bounding Box ชี้ตำแหน่งรอยโรคได้เหมือนระบบ YOLO นอกจากนี้ ด้านสถาปัตยกรรมระบบยังรันแบบต้นแบบ ทำให้กินทรัพยากรสูง และยังขาดระบบแดชบอร์ดสำหรับให้แพทย์ส่งข้อมูลย้อนกลับ (Feedback Loop)
\end{itemize}

\subsection{งานวิจัยการตรวจคัดกรองมะเร็งช่องปากแบบหลายโมดาลิตีในระดับสากล}
\textbf{Devindi et al. (IEEE Access, 2024, เอกสารเลขที่ 10664507) \cite{devindi2024multimodal}} ได้นำเสนอแนวคิด \textit{``Multimodal Deep Convolutional Neural Network Pipeline for AI-Assisted Early Detection of Oral Cancer''} ซึ่งสอดคล้องกับแนวคิดของโครงงานนี้ โดยมุ่งเน้นการบูรณาการข้อมูลภาพถ่ายจากสมาร์ตโฟนเข้ากับข้อมูลอภิพันธุ์ของผู้ป่วย (Patient Metadata เช่น ข้อมูลประวัติและอาการทางคลินิก) เพื่อเลียนแบบกระบวนการตรวจวินิจฉัยจริงของแพทย์ งานวิจัยได้เปรียบเทียบสถาปัตยกรรมตัวเข้ารหัสภาพ 6 แบบ (MobileNetV3-Large, MixNet-S, ResNet-50, HRNet-W18-C, DenseNet-121 และ Inception\_v3) ผลการวิจัยยืนยันว่าการใช้ข้อมูลภาพถ่ายร่วมกับประวัติผู้ป่วยให้ค่า Accuracy ร้อยละ 81 และ F1-score ร้อยละ 78 ซึ่งเหนือกว่าการใช้ภาพถ่ายเพียงอย่างเดียวอย่างมีนัยสำคัญ งานวิจัยนี้เป็นหลักฐานเชิงประจักษ์ระดับนานาชาติที่ยืนยันว่าการรวมข้อมูลภาพและประวัติคือทิศทางที่ถูกต้องในการคัดกรองมะเร็งช่องปาก

\subsection{การทบทวนวรรณกรรมอย่างเป็นระบบด้านการเรียนรู้เชิงลึกในมะเร็งช่องปาก}
\textbf{Warin และ Suebnukarn (BMC Oral Health, 2024, DOI: 10.1186/s12903-024-03993-5) \cite{warin2024deep}} ได้ทำการศึกษาทบทวนวรรณกรรมอย่างเป็นระบบตามแนวทางมาตรฐาน PRISMA ครอบคลุมงานวิจัยด้าน Deep Learning ในมะเร็งช่องปากทั่วโลก ผลการศึกษาชี้ให้เห็นว่าสถาปัตยกรรม CNN และ YOLO มีประสิทธิภาพสูงในการคัดกรอง แต่จุดอ่อนสำคัญของงานวิจัยในปัจจุบันคือการขาดการนำไปใช้งานจริงในคลินิก (Clinical Translation) เนื่องจากขาดระบบเชื่อมต่อที่ง่าย ขาดความสามารถในการอธิบายผล และที่สำคัญคือขาดการผสานข้อมูลบริบททางคลินิกของผู้ป่วย ซึ่งข้อเสนอแนะของงานวิจัยนี้ตรงกับวัตถุประสงค์ในการสร้างแพลตฟอร์มสนับสนุนการตัดสินใจของโครงงานเรา

\begin{table}[H]
\centering
\small
\caption{ตารางสังเคราะห์เปรียบเทียบวรรณกรรมและปริญญานิพนธ์ที่เกี่ยวข้อง}
\label{tab:synthesis_matrix}
\begin{tabular}{p{2.8cm}p{2.2cm}p{2.8cm}p{2.2cm}p{3.5cm}}
\toprule
\textbf{งานวิจัย/ปริญญานิพนธ์} & \textbf{ข้อมูลนำเข้า} & \textbf{สถาปัตยกรรมโมเดล} & \textbf{รอยโรคเป้าหมาย} & \textbf{ผลลัพธ์และข้อสังเกต} \\
\midrule
เปี่ยมเดช และ ปุณณวิช (2567) \cite{piamdech2024oral} & ภาพถ่าย + กฎเงื่อนไข & YOLOv8/v9/v10 + Rule-based & 5 รอยโรค & ตีกรอบ Bounding Box ได้ แต่ระบบแยกส่วนและเสี่ยง No Detection \\
รตา และ วรดา (2568) \cite{unnahasarakhon2025multimodal} & ภาพ + ประวัติไทย & DenseNet-121 + WangchanBERTa & 5 รอยโรค & แม่นยำ 95.5\% (Late Fusion) แต่ไม่มี Bounding Box และเป็นงานต้นแบบ \\
Devindi et al. (2024) \cite{devindi2024multimodal} & ภาพ + Metadata & MobileNetV3, DenseNet, Grad-CAM & Benign / Malignant & Accuracy 81\%, ยืนยันว่า Multimodal เหนือกว่าภาพเดี่ยว \\
Warin \& Suebnukarn (2024) \cite{warin2024deep} & การทบทวนวรรณกรรม & Systematic Review (PRISMA) & OSCC, OPMD & ชี้ชัดถึงความจำเป็นของ Multimodal, ระบบจริง และ Explainable AI \\
\textbf{โครงงานนี้ (2569)} & \textbf{ภาพ + ประวัติไทย + Feedback แพทย์} & \textbf{Hybrid: YOLO + DenseNet-BERT} & \textbf{5 รอยโรค + มะเร็งช่องปาก} & \textbf{Production Platform: มีทั้ง Bounding Box, Multimodal Fusion, Doctor Dashboard, CI/CD, PDPA} \\
\bottomrule
\end{tabular}
\end{table}

% ------------------------------------------------------------------------------
% 2.7 สรุปแนวทางการออกแบบสถาปัตยกรรมและกรอบแนวคิดของโครงงาน
% ------------------------------------------------------------------------------
\section{สรุปแนวทางการออกแบบสถาปัตยกรรมและกรอบแนวคิดของโครงงาน}
\label{sec:system_framework_summary}

จากการทบทวนวรรณกรรมและข้อจำกัดของโครงงานรุ่นพี่ทั้ง 2 รุ่น คณะผู้จัดทำจึงได้กำหนดกรอบแนวคิดในการพัฒนาระบบของโครงงานนี้ ให้ก้าวข้ามจากระดับงานวิจัยต้นแบบไปสู่ แพลตฟอร์มสนับสนุนการวินิจฉัยรอยโรคทางช่องปากที่พร้อมใช้งานจริงสำหรับทันตแพทย์ (Production-Ready Clinical Decision Support Platform) โดยมีหลักการสำคัญดังนี้:

\begin{enumerate}
    \item \textbf{การบูรณาการแบบจำลองปัญญาประดิษฐ์แบบทางเลือก (Alternative AI Model Integration):} นำแบบจำลอง YOLO สำหรับการตรวจจับและระบุตำแหน่งรอยโรค (Bounding Box) และแบบจำลองแบบหลายโหมด (Multimodal Model) มาประยุกต์ใช้งานร่วมกันภายใต้สถาปัตยกรรมไมโครเซอร์วิส (Microservices Architecture) เพื่อเพิ่มความยืดหยุ่นให้ผู้ใช้งานสามารถเลือกใช้แบบจำลองได้ตามความเหมาะสม และปรับปรุง User Interface ให้ทันสมัยและใช้งานง่ายมากขึ้น
    \item \textbf{การผสมผสานสองระบบปัญญาประดิษฐ์ (Hybrid Inference Engine):} บูรณาการทั้งโมเดล YOLO สำหรับการระบุตำแหน่งรอยโรคด้วย Bounding Box และโมเดล Multimodal สำหรับการประเมินคะแนนความเสี่ยงมะเร็งที่แม่นยำสูงเข้าด้วยกัน
    
    \item \textbf{การสร้างระบบวงจรข้อมูลย้อนกลับ (Doctor Feedback Loop):} สร้างแดชบอร์ดให้ทันตแพทย์สามารถตรวจสอบและแก้ไขผลการทำนายของ AI ได้ เพื่อเปลี่ยนแพลตฟอร์มให้กลายเป็นแหล่งรวบรวมข้อมูล Ground Truth ทางคลินิกคุณภาพสูงสำหรับการพัฒนาโมเดลต่อไป
    \item \textbf{การออกแบบสถาปัตยกรรมไมโครเซอร์วิสและ CI/CD:} แยกส่วน API Gateway (FastAPI) และ AI Engine ผ่าน Docker Containers รองรับการขยายตัวอัตโนมัติบนคลาวด์ พร้อมทั้งมีระบบ Automated Testing และ Continuous Deployment เพื่อความเสถียรสูงสุด
    \item \textbf{การรักษาความปลอดภัยและคุ้มครองข้อมูลส่วนบุคคล (PDPA Compliance):} วางระบบยืนยันตัวตน การกำหนดสิทธิ์ตามบทบาท (RBAC) และการเข้ารหัสข้อมูล เพื่อให้การประมวลผลข้อมูลสุขภาพของผู้ป่วยเป็นไปตามกฎหมายอย่างเคร่งครัด
\end{enumerate}


